%
% Capítulo 4
%
\chapter{Implementação do Sistema} \label{cap:implementacao}

A implementação do sistema foi realizada de forma incremental, começando pela definição dos fluxos principais da aplicação
Android e evoluindo depois para a integração com o backend e com a base de dados remota.
Esta abordagem permitiu validar gradualmente os principais casos de uso do projeto, garantindo que cada funcionalidade
essencial ficava operacional antes da introdução de novos componentes ou maior complexidade.

A solução resultante permite concretizar um conjunto alargado de operações centrais ao domínio do problema, nomeadamente registo e autenticação de utilizadores, definição de disponibilidades, gestão de restrições, associação entre alunos e professores, criação de propostas de horário, persistência de aulas concretas com datas reais, gestão de presenças, cancelamentos, remarcações e envio de notificações.

\section{Organização Funcional da Aplicação}

A aplicação foi concebida em torno de dois perfis principais de utilizador: professor e aluno.
A interface e os fluxos disponíveis são ajustados a cada um destes perfis, permitindo apresentar apenas as operações relevantes para cada caso.

A implementação do frontend procurou manter uma separação clara entre:

\begin{itemize}
    \item apresentação visual da interface;
    \item estado da aplicação e lógica de interação;
    \item acesso a dados locais e remotos.
\end{itemize}

Desta forma, cada ecrã da aplicação está associado a um conjunto de operações específicas,
suportadas por \textit{ViewModels} e repositórios que coordenam a lógica necessária à execução dos respetivos fluxos.

\section{Gestão de Utilizadores e Perfis}

Uma das primeiras preocupações da implementação foi distinguir corretamente os dois tipos de utilizador suportados pelo sistema.
Esta separação é essencial porque professores e alunos desempenham papéis diferentes no problema do planeamento de horários.

Na solução implementada:

\begin{itemize}
    \item o professor pode definir a sua disponibilidade, configurar restrições, consultar os alunos associados, criar propostas de horário, confirmar essas propostas e gerir as aulas persistidas;
    \item o aluno pode escolher um professor, consultar o professor associado, registar a sua disponibilidade, definir a sua carga horária semanal pretendida e consultar as aulas que lhe estão atribuídas;
    \item a navegação da aplicação adapta-se ao perfil autenticado, apresentando ecrãs e operações distintas para cada perfil.
\item \end{itemize}

A gestão de sessão local foi integrada no frontend, permitindo preservar informação sobre o utilizador autenticado e suportar a adaptação dos fluxos
da aplicação ao papel desempenhado. Esta distinção entre perfis foi também refletida no modelo de dados e na lógica dos ecrãs principais.

\section{Gestão de Disponibilidades}

A funcionalidade de gestão de disponibilidades representa uma parte fundamental da solução, pois constitui a base sobre a qual são construídas as propostas de horário.
Tanto professores como alunos podem registar a sua disponibilidade semanal, indicando os dias e intervalos horários em que se encontram disponíveis.

Numa fase inicial da implementação, a disponibilidade foi modelada de forma mais simples, com um único intervalo associado a cada dia.
Contudo, verificou-se que esta abordagem não representava adequadamente situações reais, em que um mesmo dia pode incluir múltiplos blocos horários separados.
Por esse motivo, a funcionalidade foi evoluída para suportar múltiplos intervalos por dia, permitindo representar cenários como, por exemplo,
disponibilidade durante a manhã e novamente durante a tarde no mesmo dia da semana.

Esta evolução teve impacto direto na interface e na lógica da aplicação, permitindo:

\begin{itemize}
    \item acrescentar vários intervalos no mesmo dia;
    \item definir horas de início e fim específicas para cada bloco;
    \item remover intervalos de forma individual;
    \item persistir e recuperar corretamente essa informação.
\end{itemize}


A gestão de disponibilidades foi implementada de forma a suportar registo estruturado no frontend e persistência remota no backend, mantendo no dispositivo apenas mecanismos auxiliares de suporte local.

\section{Gestão de Restrições}

Para além das disponibilidades, o sistema implementa um conjunto de restrições que orientam a construção do horário do professor.
Estas restrições representam regras de negócio relevantes para o problema e são definidas a partir do perfil do professor.

Na implementação atual, o professor pode configurar:

\begin{itemize}
    \item número máximo de horas diárias;
    \item duração de cada sessão;
    \item número máximo de alunos por sessão;
    \item número máximo de sessões por aluno por dia.
\end{itemize}

A introdução desta funcionalidade permitiu passar de uma simples compatibilização de horários para um problema de planeamento mais próximo do caso de uso real.
As restrições são persistidas e utilizadas posteriormente pelo mecanismo de criação de horários, desempenhando um papel central na validação das sessões propostas.
Para além das restrições configuradas pelo professor, o sistema passou também a permitir ao aluno definir a sua carga horária semanal pretendida. Esta informação é integrada no algoritmo de criação de horários, evitando atribuições semanais acima do limite definido para cada estudante.

\section{Associação entre Alunos e Professores}

Outro aspeto relevante da implementação foi a definição do modelo de associação entre alunos e professores.
A solução adotada considera que cada aluno pode escolher um único professor, enquanto um professor pode ter vários alunos associados.

Esta decisão foi refletida:

\begin{itemize}
    \item no modelo de dados;
    \item na lógica do frontend;
    \item nos endpoints do backend;
    \item nos ecrãs de seleção e visualização de alunos.
\end{itemize}

\section{Criação de Propostas de Horário}

A criação de horários constitui o núcleo funcional do projeto. Na versão final, este processo evoluiu de um modelo puramente abstrato para uma abordagem em dois níveis: a criação de uma proposta semanal baseada em \textit{templates} e a subsequente instanciação de aulas concretas e calendarizadas na base de dados com datas do ano civil.

A lógica de criação parte de uma representação em blocos temporais, obtidos a partir dos intervalos de disponibilidade registados. Estes blocos são analisados pelo algoritmo para identificar combinações viáveis entre o professor e os seus alunos. O processo de construção das propostas considera, de forma combinada, os seguintes aspetos:

\begin{itemize}
    \item compatibilidade de disponibilidade entre professor e alunos;
    \item limite máximo de alunos por sessão;
    \item limite de sessões por aluno no mesmo dia;
    \item limite de carga diária do professor;
    \item restrições semanais de carga horária do aluno.
\end{itemize}

Após a proposta ser validada e aceite, o sistema procede à criação em lote das aulas reais. Para gerir o vínculo e a consistência destas sessões ao longo do tempo, introduziu-se o conceito de identificador de série (\textit{seriesId}) e regras de recorrência (\textit{RecurrenceType}). Uma aula pode ser criada como um evento isolado ou fazer parte de uma sequência repetitiva (tipicamente semanalmente), sendo replicada na base de dados com datas concretas de acordo com o número estipulado de ocorrências.

Esta distinção veio conferir ao sistema a flexibilidade necessária para o quotidiano letivo, permitindo que alterações pontuais ou cancelamentos de sessões individuais coexistam com a manutenção e integridade da série de horários planeada originalmente.

\section{Visualização do Horário Criado}

Após a criação das sessões, a aplicação apresenta o horário de forma organizada e legível.
Em vez de mostrar apenas uma lista simples de resultados, foi construída uma visualização semanal com organização por dia,
onde cada sessão aparece com o respetivo intervalo horário e a lista de alunos atribuídos.

Esta representação permite ao utilizador compreender facilmente:

\begin{itemize}
    \item em que dias existem sessões;
    \item quais os intervalos temporais ocupados;
    \item que alunos foram alocados a cada sessão;
    \item em que dias não existem sessões previstas.
\end{itemize}

A implementação desta visualização constitui uma componente importante, pois transforma o resultado do algoritmo
numa representação diretamente compreensível e utilizável pelo utilizador final.
Para além da proposta semanal abstrata, a aplicação passou também a suportar a visualização de aulas concretas já persistidas, organizadas por semana e por intervalo temporal. Esta distinção entre proposta e aulas reais tornou-se importante para separar a fase de planeamento da fase operacional do sistema.

\section{Gestão Dinâmica de Aulas e Presenças}

Com o backend e as aplicações móveis totalmente integrados, foi possível expandir o sistema com um conjunto de ecrãs e fluxos operacionais destinados ao dia a dia do corpo docente e discente.

No lado do docente, através do \textit{TeacherLessonsScreen}, passou a ser possível gerir o ciclo de vida das aulas. O professor dispõe de uma vista semanal organizada por datas concretas, onde pode criar novas aulas calendarizadas, editar o horário/data (com validação automática no backend para evitar sobreposição de horários do docente) e efetuar o cancelamento individual ou em série de aulas (o que cria os referidos emails automáticos). É também neste ecrã que se faz a marcação e alteração de presenças dos alunos, cujos dados agregados alimentam um resumo estatístico consultável pelo docente.

Do lado do estudante, o ecrã \textit{StudentLessonsScreen} oferece uma visão clara e cronológica do seu horário pessoal para a semana corrente. O aluno consegue verificar o estado de cada aula (confirmada ou cancelada) e acompanhar o estado das suas aulas e a informação de assiduidade registada no sistema.

\section{Integração Progressiva com o Backend}

Ao longo da implementação, foi também sendo reforçada a integração entre o frontend e o backend.
Algumas operações começaram por ser validadas localmente, recorrendo à persistência em Room, e foram posteriormente adaptadas para consumir
a API REST e interagir com o PostgreSQL.

Este processo de integração foi realizado de forma progressiva, permitindo validar cada funcionalidade antes da passagem completa para a componente remota.
A implementação atual representa uma base funcional robusta, marcada por uma redução gradual da dependência de mecanismos locais auxiliares e pelo reforço do papel do backend como componente central do sistema.